\documentclass[10pt,a4paper]{amsart}
\usepackage[margin=19mm]{geometry}
\usepackage{amsmath,amssymb,mathtools}
\usepackage[scr=rsfs,cal=euler]{mathalfa}
\usepackage{xfrac}
\usepackage[unicode,hidelinks,bookmarks=false]{hyperref}
\newcommand{\ie}{i.e.\ }
\newcommand{\cf}{cf.\ }
\newcommand{\resp}{respectively\ }
\newcommand{\Subsection}{\subsection}
\providecommand{\nobreakdash}{\nobreak\hbox{-}}
\setlength{\emergencystretch}{2em}
\multlinegap0pt
\usepackage{amssymb}
\usepackage{tikz}
\newtheorem{lemma}{Lemma}
\newcommand*{\Ge}{\geqslant}
\newcommand*{\hh}{\mathcal{H}}
\newcommand*{\nul}{\mathscr{N}}
\newcommand*{\ogr}{\boldsymbol{B}}
\newcommand*{\ran}{\mathscr{R}}
\newcommand*{\zbb}{\mathbb{Z}}
\title{Hyperrigidity I: singly generated commutative $C^*$-algebras}
\author{Pawe{\l} Pietrzycki, Jan Stochel}
\date{Source: arXiv:2405.20814v4 (CC BY 4.0)}
\begin{document}
\maketitle

\noindent\emph{Source and licence.} Hyperrigidity I: singly generated commutative $C^*$-algebras. Version arXiv:2405.20814v4 (CC BY 4.0), distributed by arXiv under the Creative Commons Attribution 4.0 International licence, http://creativecommons.org/licenses/by/4.0/, as recorded in the arXiv licence metadata for this exact version and shown on the abstract page https://arxiv.org/abs/2405.20814v4. Full licence notice: \texttt{licenses/arxiv-2405.20814-CC-BY-4.0.txt}.\par\smallskip

\noindent\textit{Editorial scope.} Counted unit: the article's lemma npq (source0.tex lines 823--847). The article's complete original proof of it is delivered with it (source0.tex lines 848--907, one \texttt{proof} environment of 60 lines). No mathematics is written, repaired or re-derived: the statement and the proof are the article's own lines, in the article's own notation, and no retained line is substituted or re-flowed.  Retained uncounted as the source-local context the proof needs: lines 624--627; lines 634--637.  Nothing was omitted from any retained span, so the package carries no omission record.  No retained line contains a citation, so the wrapper carries no bibliography and no bibliography entry.  No retained line contains a figure environment, an image-inclusion command or drawing code, so no image is delivered and the package declares no asset dependency.  Editorial formatting only, in addition: an AMS \texttt{amsart} preamble that restates the article's own declarations and macros used by the retained lines, namely the environments \texttt{lemma} on separate counters, together with the standard packages those lines need.  The retained environments print in this document as Lemma 1 in order of appearance, whereas the article prints the same items with the numbering of its own full document.  The complete native source of the article as distributed in the arXiv e-print for this exact version is bundled unchanged as \texttt{sources/160065/source0.tex}.
   \begin{align} \label{pdu}
\text{$N=U |N|=|N|U$, \, $U$ is a unitary operator \,
and \, $U|_{\nul(N)}=I_{\nul(N)}$.}
   \end{align}

   \begin{align} \label{orotrh}
V=\tilde V \oplus 0, \quad U=\tilde V \oplus
I_{\nul(N)},
   \end{align}

   \begin{lemma} \label{npq}
Let $N\in \ogr(\hh)$ be a normal operator,
$N=U |N|$ be the unitary polar decomposition
of $N$ $($see \eqref{pdu}$)$, $P \in
\ogr(\hh)$ be an orthogonal projection and
$(p,q)\in \zbb_+^2 \setminus \{(0,0)\}$.
Then the following conditions are
equivalent{\em :}
   \begin{enumerate}
   \item[(i)] $P$ commutes with $N^{*p}N^q$,
   \item[(ii)] $P$ commutes with $|N|$ and
$U^{q-p}$.
   \end{enumerate}
Moreover, if $p\neq q$, then {\em (i)} is
equivalent to any of the following
statements{\em :}
   \begin{enumerate}
   \item[(iii)] $P$ commutes with
$N^{*{\tilde p}}N^{\tilde q}$ for every
$(\tilde p, \tilde q) \in \zbb_+^2 \setminus
\{(0,0)\}$ such that $|\tilde q-\tilde
p|=|q-p|$,
   \item[(iv)] $P$ commutes with $N^{|q-p|}$.
   \end{enumerate}
   \end{lemma}

   \begin{proof}
Since $U$ and $|N|$ commute (see
\eqref{pdu}), we see that
   \begin{align} \label{pulder}
N^{*p}N^q=U^{q-p}|N|^{p+q}.
   \end{align}

(i)$\Rightarrow$(ii) First, we show that $P$ and $|N|$
commute. It follows from the selfadjointness of $P$
that $P$ commutes with $(N^{*p}N^q)^*$, and therefore
also with
   \begin{equation*}
(N^{*p}N^q)^*N^{*p}N^q=|N|^{2(p+q)}.
   \end{equation*}
Since $p+q \Ge 1$, we see that the commutants of
$|N|^{2(p+q)}$ and $|N|$ coincide, which yields
$P|N|=|N|P$. This, together with (i) and
\eqref{pulder}, leads to
   \begin{equation} \label{upn}
PU^{q-p}|N|^{p+q}=U^{q-p}|N|^{p+q}P=U^{q-p}P|N|^{p+q}.
   \end{equation}
The kernel-range decomposition, together with $p+q \Ge
1$, yields
   \begin{equation*}
\overline{\ran(|N|^{p+q})}^\perp =
\nul(|N|^{p+q})=\nul(|N|) =
\overline{\ran(|N|)}^\perp,
   \end{equation*}
which implies that
$\overline{\ran(|N|^{p+q})} =
\overline{\ran(|N|)}$. This, combined with
\eqref{upn}, shows that
   \begin{equation}\label{pur}
PU^{q-p}|_{\overline{\ran(|N|)}}=U^{q-p}P|_{\overline{\ran(|N|)}}.
   \end{equation}
On the other hand, the fact that $P$
commutes with $|N|$ implies that
$P\nul(|N|)\subseteq \nul(|N|)$. Since by
\eqref{orotrh},
$U|_{{\nul(|N|)}}=I_{{\nul(|N|)}}$, we get
   \begin{equation*}
PU^{q-p}|_{{\nul(|N|)}}=P|_{{\nul(|N|)}}=U^{q-p}P|_{{\nul(|N|)}}.
   \end{equation*}
This combined with \eqref{pur} shows that
$P$ commutes with $U^{q-p}$.

(ii)$\Rightarrow$(i) Use \eqref{pulder}.

Assume now that $p\neq q$.

(ii)$\Rightarrow$(iii) Apply \eqref{pulder} to
$(\tilde p, \tilde q)$ in place of $(p,q)$.

(iii)$\Rightarrow$(iv) Apply (iii) to
$\tilde p= 0$ and $\tilde q=|q-p|$.

(iv)$\Rightarrow$(ii) Apply implication
(i)$\Rightarrow$(ii) to the pair $(0,|q-p|)$
in place of $(p,q)$.
   \end{proof}

\end{document}
